\section*{Capítulo \ref{ch:g8:the-oceans} --- Los océanos}

\begin{solution}{exo:g8:the-oceans:1}
\omterm{def:g8:the-oceans:seafloor}{Plataforma continental}, \omterm{def:g8:the-oceans:seafloor}{talud continental}, \omterm{def:g8:the-oceans:seafloor}{llanura abisal}, \omterm{def:g8:the-oceans:seafloor}{dorsal oceánica}.
\end{solution}

\begin{solution}{exo:g8:the-oceans:2}
(a) Falso: unos \qty{3700}{m}, es decir, \qty{3.7}{km}. (b) Verdadero.
(c) Falso: la evaporación quita agua y deja las sales, así que el agua se
vuelve más salada. (d) Verdadero.
\end{solution}

\begin{solution}{exo:g8:the-oceans:3}
(a) Una \omterm{def:g8:the-oceans:seafloor}{dorsal oceánica}. (b) La \omterm{def:g8:the-oceans:salinity}{salinidad}. (c) Un \omterm{def:g8:the-oceans:current}{giro oceánico}. (d) Una \omterm{def:g8:the-oceans:seafloor}{fosa oceánica}.
\end{solution}

\begin{solution}{exo:g8:the-oceans:4}
(a) En sentido contrario a las agujas del reloj. (b) Cálida: la corriente
del Golfo; fría: la corriente de Canarias. (c) Los \omterm{def:g8:winds-global-circulation:global-winds}{alisios}.
\end{solution}

\begin{solution}{exo:g8:the-oceans:5}
(a) $35 \div 1000 \times 100 = 3.5~\%$. (b) $20 \times 35 = \qty{700}{g}$.
(c) $1000 \div 35 \approx 28.6$: unos \qty{29}{kg}.
\end{solution}

\begin{solution}{exo:g8:the-oceans:6}
(a) $1500 \times 6 = \qty{9000}{m}$. (b) El sonido baja hasta el fondo y
vuelve a subir: recorre dos veces la profundidad. (c) $9000 \div 2 =
\qty{4500}{m}$. (d) Una \omterm{def:g8:the-oceans:seafloor}{llanura abisal} (una fosa sería más profunda, y una
dorsal o una plataforma, menos).
\end{solution}

\begin{solution}{exo:g8:the-oceans:7}
(a) Entre unos \qty{3500}{m} y \qty{5000}{m}. (b) A unos \qty{4400}{km} a
lo largo del perfil, a unos \qty{2300}{m} de profundidad. (c) Unos
$4500 - 2300 = \qty{2200}{m}$ (entre 1200 y \qty{2700}{m}, según la
\omterm{def:g1:landscapes-around-us:landform}{llanura} elegida). (d) El \omterm{def:g2:water-on-earth:ocean}{océano} tiene unos pocos kilómetros de profundidad
pero miles de kilómetros de anchura: dibujado a la misma escala, todo el
perfil sería más fino que una línea.
\end{solution}

\begin{solution}{exo:g8:the-oceans:8}
El \omterm{def:g2:water-on-earth:ocean}{mar} A es más salado: el Sol y el aire seco provocan una fuerte
\emph{evaporación}, el agua se va, \emph{las sales se quedan}, y casi
ninguna \omterm{def:g4:clouds-rain-snow:precipitation}{lluvia} ni ningún \omterm{def:g1:landscapes-around-us:water}{río} devuelve \emph{\omterm{def:g2:water-on-earth:salt-fresh}{agua dulce}}. El \omterm{def:g2:water-on-earth:ocean}{mar} B es menos
salado: sus \omterm{def:g1:landscapes-around-us:water}{ríos} y su \omterm{def:g4:clouds-rain-snow:precipitation}{lluvia} aportan mucha \emph{\omterm{def:g2:water-on-earth:salt-fresh}{agua dulce}}, que diluye
las sales.
\end{solution}

\begin{solution}{exo:g8:the-oceans:9}
(a) La corriente del Golfo, hacia el norte a lo largo de la \omterm{def:g1:landscapes-around-us:water}{costa}, y
después hacia el nordeste. (b) Hacia el este (el nordeste), con la
corriente del Atlántico Norte. (c) A lo largo de la \omterm{def:g1:landscapes-around-us:water}{costa} del noroeste de
África (y de España y Portugal), con la corriente de Canarias. (d) Fría.
\end{solution}

\begin{solution}{exo:g8:the-oceans:10}
(a) $6.4 \times 24 = 153.6$: unos \qty{154}{km}. (b) $153.6 \times 7
\approx \qty{1075}{km}$. (c) La corriente del Golfo es un poco más rápida
que un caminante, pero nunca se detiene, ni de día ni de noche.
\end{solution}

\begin{solution}{exo:g8:the-oceans:11}
(a) \qty{3.7}{km} y \qty{11}{km}. (b) $11000 \div 3700 \approx 3.0$
veces. (c) El sonido recorre $2 \times 11000 = \qty{22000}{m}$; $22000
\div 1500 \approx 14.7$: unos \qty{15}{s}. (d) $1000 = 10^3$ años.
\end{solution}

\begin{solution}{exo:g8:the-oceans:12}
(a) Cerca del \omterm{def:g3:continents-and-oceans:equator}{ecuador}, los \omterm{def:g8:winds-global-circulation:global-winds}{alisios} empujan el agua tropical cálida hacia el
oeste. En el borde occidental del Atlántico gira hacia el norte con la
corriente del Golfo, una corriente cálida, a lo largo de Florida y de
América del Norte. Después, los \omterm{def:g8:winds-global-circulation:global-winds}{vientos del oeste} la empujan a través del
Atlántico hacia el nordeste, hasta los \omterm{def:g2:water-on-earth:ocean}{mares} de Noruega. (b) Allí cede su
calor al aire frío: se enfría, y se vuelve más salada cuando se forma el
hielo marino, porque la sal se queda en el agua. El agua fría y salada es
densa, así que se hunde hasta el fondo. (c) Avanza lentamente hacia el sur
a ras del fondo oceánico, hacia la Antártida, y después hacia los \omterm{def:g2:water-on-earth:ocean}{océanos}
Índico y Pacífico, donde vuelve a subir poco a poco.
\end{solution}

\begin{solution}{pb:g8:the-oceans:1}
\textbf{1.} En el \omterm{def:g2:water-on-earth:ocean}{océano} Atlántico, en el \omterm{def:g3:continents-and-oceans:equator}{hemisferio norte}.

\textbf{2.} En el sentido de las agujas del reloj; la boya es arrastrada primero hacia el norte.

\textbf{3.} El tramo A (corriente del Golfo) es cálido: viene de los
\omterm{def:g5:weather-and-climate:tropics}{trópicos}. El tramo C (corriente de Canarias) es frío: viene de \omterm{def:g6:latitude-longitude:latitude}{latitudes}
más altas, hacia el \omterm{def:g3:continents-and-oceans:equator}{ecuador}.

\textbf{4.} Tramo B: los \omterm{def:g8:winds-global-circulation:global-winds}{vientos del oeste}. Tramo D: los \omterm{def:g8:winds-global-circulation:global-winds}{alisios}.

\textbf{5.} $3800 \div 100 = 38$ días.

\textbf{6.} B: $2100 \div 20 = 105$ días. C: $3100 \div 20 = 155$ días.
D: $6100 \div 20 = 305$ días.

\textbf{7.} $3800 + 2100 + 3100 + 6100 = \qty{15100}{km}$.

\textbf{8.} La corriente del Golfo es la corriente estrecha y rápida del
lado occidental del \omterm{def:g8:the-oceans:current}{giro}; las corrientes de los otros lados son anchas y
lentas.

\textbf{9.} $38 + 105 + 155 + 305 = 603$ días.

\textbf{10.} $603 \div 30 = 20.1$: unos \textbf{20 meses}.

\textbf{11.} Después de los tramos A, B y C han pasado
$38 + 105 + 155 = 298$ días; los 365 días llegan durante el tramo D (la
corriente norecuatorial), $365 -
298 = 67$ días después de empezarlo.

\textbf{12.} $603 - 365 = 238$ días más, es decir, $238 \div 30 \approx 8$
meses más.
\end{solution}
